\chapter*{Abstract}
\addcontentsline{toc}{chapter}{Abstract}

Simulating a soft polymer that carries rigid particles normally requires a
microstructural finite-element (FE) problem at every macroscopic integration
point. This Personal Programming Project implements, verifies and critically
assesses an offline alternative for polydimethylsiloxane (PDMS) reinforced by
rigid glass spheres. A periodic representative volume element (RVE) with a
near-incompressible mixed $P_2/P_1$ matrix and an \emph{exactly} rigid,
freely rotating particle is solved in NGSolve. Its work-conjugate homogenized
stress follows from the assembled particle reaction, and its consistent
$9\times9$ tangent from Schur condensation of the discrete energy Hessian. The
resulting labels train a volume-fraction-conditioned, partially input-convex
neural network (pICNN) in PyTorch, whose energy is differentiated
automatically and coupled to a total-Lagrangian NGSolve Newton solver.

All 26 verification tests pass. These include the guideline-mandated checks:
six rigid-body zero eigenvalues, the twelve hourglass modes of an under-integrated hexahedron, and
single-element tests under full deformation and under stress control. The RVE satisfies periodicity and rigidity to $10^{-16}$ and the
Hill--Mandel condition to $10^{-15}$, and its tangent shows Taylor remainder
order $2.000$. Its small-strain moduli agree with the Mori--Tanaka rigid-sphere
estimate within $1.1\%$.

The surrogate assessment is deliberately split into pressure and deviatoric
stress. The pICNN reproduces pressure to a median $0.03\%$ but deviatoric stress
only to a median $76$--$83\%$ outside expansion states, and it overestimates the small-strain shear
modulus by $46$--$70\%$. The cause is identified quantitatively: with
$\kappa/\mu=10^4$ the training loss is dominated by volumetric energy, so that
removing the entire deviatoric stress would raise the validation stress loss
by only $1.5\%$. The
frozen model replaces a $134\,$s RVE solve by a $2.3\,$ms evaluation, a median
speed-up of about $5.8\times10^4$. It runs inside a macroscopic FE solve to
$5\%$ stretch with linear elements, although this solution is shown to be
volumetrically locked, while quadratic elements reach $1\%$ on two of three meshes. The project
therefore delivers a verified multiscale software pipeline, a diagnosed
limitation of the surrogate's accuracy, and a concrete remedy: isochoric
polyconvex features with a split loss, and a mixed macroscopic formulation.

\medskip\noindent\textbf{Keywords:} computational homogenization;
hyperelasticity; input-convex neural network; representative volume element;
NGSolve; PyTorch; PDMS; verification.
